\begin{textsample}{POS Dim 1 – vem\_ed\_33 – Score 15.00 – vem\_ed\_33/t176.txt}  \label{ex:f1_pos_078}
18ª Semana de Práticas e Atualizações Pedagógicas A equipe de Apoio à Internacionalização do Ensino Superior (AIES) da Divisão de Extensão e Pesquisa (Depes) participou da organização da 18ª Semana de Práticas e Atualizações Pedagógicas \textbf{promovida} entre os dias 4 e 6 de fevereiro de 2026 pela Coordenadoria de Ensino Superior de Graduação (CGESG) do Centro Paula Souza (CPS).
Patrícia Patrício, responsável pela Comunicação da equipe, \textbf{foi} mestre de cerimônias.
Maria Claudia Nunes Delfino fez a curadoria das perguntas enviadas pelos participantes.
Neusa Haruka Sezaki Gritti e Priscilla Ferro ficaram nos bastidores acompanhando a \textbf{interação} entre o público e os palestrantes.
No dia 6 de fevereiro, Osvaldo Succi Junior palestrou sobre os Projetos Colaborativos Internacionais (PCIs) como atividade de \textbf{extensão} \textbf{institucional} das Fatecs.
Succi Junior pontuou \textbf{que} o \textbf{intercâmbio} pedagógico entre instituições nacionais e estrangeiras estava previsto desde 1938 pelo Instituto Nacional de Estudos e Pesquisas Educacionais Anísio Teixeira (Inep).
Confira a gravação em https://tinyurl.com/shbjnup3.
A Portaria da Presidência do Centro Paula Souza, \textbf{publicada} em janeiro de 2026, formaliza a promoção da internacionalização com o objetivo de estimular a presença e reputação do CPS no cenário internacional.
Nesse sentido, os PCIs, tropicalização do modelo COIL (Collaborative Online International Learning), \textbf{permitem} \textbf{interação} virtual entre professores e alunos de diferentes países e aprimoram \textbf{competências} \textbf{socioemocionais}, \textbf{linguísticas}, digitais e de trabalho em equipes multiculturais.
Para os docentes das Fatecs, abrem portas para publicações conjuntas com os colegas internacionais, \textbf{compartilhando} os resultados dos PCIs.
Desde 2013, \textbf{foram} realizados 750 PCIs envolvendo 57 Fatecs e cerca de 94 instituições estrangeiras, engajando cerca de 15 mil fatecanos e 15 mil estudantes das instituições parceiras, o \textbf{que} coloca as Fatecs do Centro Paula Souza no"Top 5"mundial em \textbf{número} de projetos COIL.
A equipe AIES participa de redes de colaboração internacional em Intercâmbios Virtuais como COIL Connect, Red LatAM COIL, UniCollaboration e Faubai (Associação Brasileira de Educação Internacional).
Em outubro de 2022, o CPS conquistou o selo BraVE (Brazilian Virtual Exchange) da Faubai, \textbf{reconhecendo} a qualidade dos PCIs.
Os temas abordados nos PCIs \textbf{são} \textbf{variados} como a oferta de cursos nas Fatecs.
Apenas para citar alguns exemplos da multiplicidade de assuntos, já \textbf{foram} realizados PCIs sobre publicidade, sustentabilidade no setor têxtil, agricultura, poluição do ar, comunicação empresarial, entre outros.
Cerca de 75 \% dos projetos \textbf{são} realizados em inglês, 23 \% em espanhol e 2 \% em português.

% matched lemmas: compartilhar, competência, extensão, institucional, interação, intercâmbio, linguístico, número, permitir, promover, publicar, que, reconhecer, ser, socioemocional, variar
\end{textsample}
